\documentclass[10pt,a4paper]{amsart}
\usepackage[margin=19mm]{geometry}
\usepackage{amsmath,amssymb,mathtools}
\usepackage[scr=rsfs,cal=euler]{mathalfa}
\usepackage{xfrac}
\usepackage[unicode,hidelinks,bookmarks=false]{hyperref}
\newcommand{\ie}{i.e.\ }
\newcommand{\cf}{cf.\ }
\newcommand{\resp}{respectively\ }
\newcommand{\Subsection}{\subsection}
\providecommand{\nobreakdash}{\nobreak\hbox{-}}
\setlength{\emergencystretch}{2em}
\multlinegap0pt
\usepackage{amssymb}
\usepackage{xcolor}
\usepackage{tikz}
\usepackage{dsfont}
\usepackage{enumerate}
\newtheorem{theorem}{Theorem}
\newtheorem{lemma}{Lemma}
\DeclareMathOperator{\bound}{\Big( 3d-1, \coe\sqrt{n}\Big)}
\DeclareMathOperator{\coe}{6\Big(1+\frac{9d^2}{4}\Big)^{\frac{1}{2}}d}
\DeclareMathOperator{\ex}{\text{Ex}}
\DeclareMathOperator{\mad}{\text{mad}}
\title{MAD Phase Transitions in the Oriented Chromatic Number}
\author{Alexander Clow\,}
\date{Source: arXiv:2607.22915v1 (CC BY 4.0)}
\begin{document}
\maketitle

\noindent\emph{Source and licence.} MAD Phase Transitions in the Oriented Chromatic Number. Version arXiv:2607.22915v1 (CC BY 4.0), distributed by arXiv under the Creative Commons Attribution 4.0 International licence, http://creativecommons.org/licenses/by/4.0/, as recorded in the arXiv licence metadata for this exact version and shown on the abstract page https://arxiv.org/abs/2607.22915v1. Full licence notice: \texttt{licenses/arxiv-2607.22915-CC-BY-4.0.txt}.\par\smallskip

\noindent\textit{Editorial scope.} Counted unit: the article's theorem Upper (source0.tex lines 324--328). The article's complete original proof of it is delivered with it (source0.tex lines 620--715, one \texttt{proof} environment of 96 lines, which the article places away from the statement). No mathematics is written, repaired or re-derived: the statement and the proof are the article's own lines, in the article's own notation, and no retained line is substituted or re-flowed.  Retained uncounted as the source-local context the proof needs: lines 347--349; lines 377--381; lines 457--464; lines 532--543.  Nothing was omitted from any retained span, so the package carries no omission record.  No retained line contains a citation, so the wrapper carries no bibliography and no bibliography entry.  No retained line contains a figure environment, an image-inclusion command or drawing code, so no image is delivered and the package declares no asset dependency.  Editorial formatting only, in addition: an AMS \texttt{amsart} preamble that restates the article's own declarations and macros used by the retained lines, namely the environments \texttt{theorem}, \texttt{lemma} on separate counters, together with the standard packages those lines need.  The retained environments print in this document as Theorem 1, Lemma 1 in order of appearance, whereas the article prints the same items with the numbering of its own full document.  The complete native source of the article as distributed in the arXiv e-print for this exact version is bundled unchanged as \texttt{sources/205983/source0.tex}.
\begin{lemma}\label{Lemma: MAD}
    If $G$ is a $d$-degenerate graph, then $\mad(G) < 2d$.
\end{lemma}

\begin{lemma}\label{Lemma: counter-no-oneway}
    If $G$ is a smallest counterexample to Theorem~\ref{Thm: Critical Upper},
    then for all vertices $v \in V(G)$ with degree at most $3d-1$,
    $\deg(v)^+,\deg^-(v)\geq 1$.
\end{lemma}

\begin{lemma}\label{Lemma: extendedEquiv}
    Let $T$ be a $(3d-1,1)$-comprehensive graph,
    let $G$ be an oriented graph, and let $\sigma$ be a vertex ordering of $G$.
    If the vertices of $G$ with degree at most $3d-1$ form an independent set,
    and no vertex with degree at most $3d-1$ is a source or a sink,
    then $G$ has a homomorphism to $T$ if and only if 
    $\ex_d(G,\sigma)$ has a homomorphism to $T$.
\end{lemma}

\begin{lemma}\label{Lemma: The Fun Bit}
    Let $G$ be a smallest counterexample to Theorem~\ref{Thm: Critical Upper}, let $\sigma$ be a vertex ordering of $G$,
    let $A$ be the set of vertices with degree at least $3d$ in $\ex_d(G,\sigma)$, and let $B$ be the set of all vertices of degree at most $2$ in $\ex_d(G)$.
    Then 
    for all vertices $v \in A$, if $v$ has $k\leq 3d-1$ neighbours in $A$, 
    then 
    there are at least
    \[
    (\coe\sqrt{n})2^{-k+(3d-1)}.
    \]
    distinct vertices $w\in A$ such that $N(v)\cap N(w)\cap B \neq \emptyset$.
\end{lemma}

\begin{theorem}\label{Thm: Critical Upper}
    Let $d\geq 2$ be an integer and let $G$ be an orientation of a $d$-degenerate graph.
    If $G$ has $n$ vertices
    and $T$ is a $\bound$-comprehensive graph, then $G$ admits an oriented homomorphism to $T$.
\end{theorem}

\begin{proof}[Proof of Theorem~\ref{Thm: Critical Upper}]

We will prove the theorem by induction of $n$.
If $n\leq 1$, then trivially the result holds.
Suppose then that $n>1$ is the least integer such that
there exists a counterexample $G$ to the theorem with $n(G) = n$.
Let $G$ be such a counterexample.
Let $T$ be a $\bound$-comprehensive graph
such that $G$ has no homomorphism to $T$.
Notice $T$ exists since $G$ is a counterexample.

As in Lemma~\ref{Lemma: extendedEquiv} and Lemma~\ref{Lemma: The Fun Bit}
we write $\ex_d(G)$ for $\ex_d(G,\sigma)$ since the choice of $\sigma$ does not matter.
Let $A$ be the set of vertices in $\ex_d(G)$ with degree greater than $(3d-1)$,
and let $B$ be the set of degree at most $2$ vertices in $\ex_d(G)$.
By Lemma~\ref{Lemma: counter-no-oneway} every vertex in $B$ has degree $2$.
Next, 
by Lemma~\ref{Lemma: The Fun Bit} for every vertex $v \in A$ with $k \leq (3d-1)$ neighbours in 
$A$
there are at least
    \[
    (\coe\sqrt{n})2^{-k+(3d-1)}.
    \]
distinct vertices $w\in A$ such that $N(v)\cap N(w)\cap B \neq \emptyset$.
We note that since $G[A]$ and $\ex_d(G)[A]$ are isomorphic,
and since $G$ is $d$-degenerate there exists a vertex in $A$ with at most $d$ neighbours in $A$.
Let $N = |A|$.

Now define $H$ from $\ex_d(G)$ by deleting vertices from $B$ until there are no twin vertices in $B$.
That is until there are no vertices $u,v \in B$ with the same in-neighbours and out-neighbours.
Trivially, $H$ has an oriented homomorphism to $T$ if and only if $\ex_d(G)$ does.
Hence, $H$ has no oriented homomorphism to $T$.
Going forward we will only consider $H$, notice that $A$ is unchanged.
We let $B'$ denote the subset of $B$ in $H$.

Since $B$ is an independent set in $\ex_d(G)$, $B'$ is an independent set in $H$.
Thus, all vertices in $B'$ are degree $2$, with their neighbours in $A$, and $B'$ contains no twins.
Hence, 
\begin{align*}
    n(H) & \leq N +2\binom{N}{2}\\
    & = N^2
\end{align*}
implying $\sqrt{n(H)} \leq N$.
Similarly, notice that each vertex $v$ of degree at most $3d-1$ which is 
removed from $G$ and replaced by degree $2$ vertices in $\ex_d(G)$, 
adds exactly $\deg^-(v)\deg^+$ degree $2$ vertices.
Elementary calculus implies $\deg^-(v)\deg^+\leq \frac{deg(v)^2}{4} < \frac{9d^2}{4}$. 
Since at most $n$ vertices are deleted and replaced with degree $2$ vertices 
during the construction of $\ex_d(G)$ we observe that
\begin{align*}
    n(H) & \leq |V(\ex_d(G))| \\
    & \leq n + \frac{9d^2}{4}n
\end{align*}
implying that $n(H) \leq \Big(1+\frac{9d^2}{4}\Big)n$.
Observe that since $G$ is $d$-degenerate, $\ex_d(G)$ is $d$-degenerate, implying $H$ is $d$-degenerate, 
since $H$ is a subgraph of $\ex_d(G)$.


Since $H$ is $d$-degenerate Lemma~\ref{Lemma: MAD} implies
$\mad(H) < 2d$. So the average degree of $H[A]$ is strictly less than $2d$.
Let $X$ and $Y$ be a partition of $A$ such that $v \in X$
if and only if $\deg_{H[A]}(v)\geq 3d$.
We observe that
\begin{align*}
    2d &> \frac{\sum_{v \in A} \deg_{H[A]}(v)}{N} \\
    & \geq \frac{3d|X|}{N}
\end{align*}
implies that $|X| < \frac{2N}{3}$. From which it follows that 
\[
|Y| > \Big(1-\frac{2}{3}\Big)N\geq  \frac{1}{3}\sqrt{n(H)}.
\]
Now
by Lemma~\ref{Lemma: The Fun Bit} and the definition of $H$ for all vertices $v \in A$ with $\deg_{H[A]}(v) \leq (3d-1)$,
\[
\deg_H(v) \geq  \deg_{H[A]}(v) + \Big( \coe \sqrt{n}\Big)2^{-\deg_{H[A]}(v)+(3d-1)}.
\]
Observe that this function is monotone decreasing for values of $\deg_{H[A]}$ at most $3d-1$.
From here we consider the average degree of the entire graph $H$,
\begin{align*}
    2d &> \frac{\sum_{v \in V(H)} \deg_{H}(v)}{n(H)} \\
    \\
    & \geq \frac{\sum_{v \in Y} \deg_{H}(v)}{n(H)} \\
    \\
    & \geq \frac{1}{n(H)} \Bigg(\sum_{v\in Y} \Big(\deg_{H[A]}(v) + \Big(\coe\sqrt{n}\Big)2^{-\deg_{H[A]}(v)+(3d-1)}\Big)\Bigg) \\
    \\
    &\geq \frac{|Y|}{n(H)}\Bigg(\coe\sqrt{n} \Bigg)  \\
    \\ 
    &>\frac{1}{3\sqrt{n(H)}}\Bigg(\coe\sqrt{n} \Bigg) \\
    \\
    & \geq \frac{\coe }{3} \Bigg(1+\frac{9d^2}{4}\Bigg)^{-\frac{1}{2}}\\
    \\
    & = 2d
\end{align*}
a contradiction. But we have made no assumption except that $G$ is a smallest counterexample.
So there is no counterexample, completing the proof.
\end{proof}

\end{document}
